\section{Training Hyperparameters}
\label{app:hyperparams}

% ⚠ SUPERSEDED -- kept commented for comparison, not deleted. Everything below described the
% earlier LoRA/Axolotl regime and contradicted the main text, which states full fine-tuning at
% lr 5e-5 (Section 4.1). The CTC-suite runs are full-parameter fine-tunes in olmo-core.
% Following initial experiments where full-finetuning seems to performan similarly or worse on our settings, we train on our settings using LoRA \citep{Hu2021LoRALA} fine-tuning on a frozen base model, AdamW-8bit, cosine LR schedule, and bfloat16. We use Axolotl \citep{axolotl} for our standard training experiments, and custom code based on the Huggingface \cite{wolf-etal-2020-transformers} transformers library.  Defaults shared across all settings following valdation from initial hyperparameter tuning:
% \begin{itemize}
% \item \textbf{LoRA:} $r{=}32$, $\alpha{=}64$, dropout 0.05, target modules $\{q,k,v,o,\text{gate},\text{up},\text{down}\}_\text{proj}$.
% \item \textbf{Optimizer:} AdamW-8bit, lr $1\!\times\!10^{-4}$, cosine schedule, weight decay 0, warmup 5--10\%.
% \item \textbf{Precision:} bf16 weights, tf32 matmul, gradient checkpointing on, no flash-attention (since chunked attention paths require a custom mask).
% \item \textbf{Packing:} \texttt{sample\_packing} is \emph{off} for all results in the main paper, so per-example sequence lengths are stable across runs.
% \item \textbf{DeepSpeed:} ZeRO-2 for multi-GPU runs.
% \end{itemize}

% \amandanote{ah, we should discuss whether full finetuning might make chunked attn seem stronger (and maybe consider using LongLoRA in an ablation for the rebuttal? I would ask for that if I reviewed this paper.)} \prasannnote{Hmm yeah I guess FFT I checked a bit earlier into the project and it was not doing so great, but LongLoRA is a really good point...}

\begin{claudeblock}
All CTC Suite results are \textbf{full-parameter fine-tunes} of a base checkpoint, trained with
\texttt{olmo-core}. Every arm of a comparison --- full attention, document-chunked, and
document-chunked with mask mixing --- is trained from the same base with the same recipe, so the
only difference between arms is the attention mask used during training. Defaults shared across
all settings:

\begin{itemize}
\item \textbf{Optimizer:} AdamW with skip-step (a step whose loss or gradient norm is a large
outlier is skipped rather than applied), lr $5\!\times\!10^{-5}$, $\beta = (0.9, 0.95)$, weight
decay 0, gradient clipping at global norm 1.0.
\item \textbf{Schedule:} linear warmup over the first 3\% of steps, then linear decay to 0.
\item \textbf{Epochs / batch:} 1 epoch, global batch 8 sequences, 1 instance per micro-batch.
\item \textbf{Precision:} bf16 parameters with fp32 gradient reduction, FSDP for data parallelism,
activation checkpointing on.
\item \textbf{Packing:} examples are \emph{not} packed; one example per sequence, so per-example
sequence lengths are stable across runs and document boundaries stay aligned with chunk boundaries.
\item \textbf{Sequence length:} set per task from the maximum example length of that task's
tokenized shard --- 40960 for contradiction, reorder, QDmatch and outlier, and 26112 for FiQA.
This matters: the data converter \emph{drops} over-length examples, so a sequence length below the
shard maximum would silently delete the long tail, which on a length ladder is exactly the part
being measured.
\end{itemize}

\texttt{torch.compile} is disabled for the mask-mixing arm, whose per-step choice of mask is not
compilable; this is a speed difference only, and the mask itself is identical to the pure-chunked
arm's.
\end{claudeblock}

\claude{Because the recipe above is shared, per-setting variation reduces to which models were
trained on a task, the ladder it was evaluated on, and the sequence length; these are summarized in
Table~\ref{tab:hyperparams}.}

% ⚠ SUPERSEDED -- the old per-task table is kept commented directly below its replacement. It
% listed the LoRA-era knobs (micro-batch / grad-accum / 3--4 epochs) and models that no longer
% appear in the paper (Qwen-3.5-9B reorder, Qwen-2.5-3B contradiction), at 4k--16k rather than the
% 2k--32k ladders actually used.

\begin{table}[htbp]
\centering
\small
\setlength{\tabcolsep}{4pt}
\begin{claudeblock}
\begin{tabular}{l l l c c}
\toprule
\textbf{Task} & \textbf{\cname} & \textbf{Models trained} & \textbf{Ladder} & \textbf{seq\_len} \\
\midrule
HotpotQA        & \ctc{N}   & Qwen3.5-\{0.8B, 2B, 4B\}, Olmo3-7B, Llama3.2-3B & 2k--32k & 40960 \\
BEIR FiQA       & \ctc{N}   & Qwen3.5-\{0.8B, 4B\}                            & 2k--32k & 26112 \\
\midrule
Outlier (scale-$k$) & \ctc{NM} & Qwen3.5-4B                                   & 2k--32k & 40960 \\
\midrule
QDmatch (NQ)    & \ctc{N^2} & Qwen3.5-\{0.8B, 2B, 4B\}                        & 2k--32k & 40960 \\
QDmatch (FiQA)  & \ctc{N^2} & Qwen3.5-4B                                      & 2k--32k & 40960 \\
Reordering      & \ctc{N^2} & Qwen3.5-\{0.8B, 2B, 4B\}                        & 2k--16k & 40960 \\
\multirow{2}{*}{Contradiction} & \multirow{2}{*}{\ctc{N^2}}
                & Qwen3.5-\{0.8B, 2B, 4B\}, Olmo3-7B,                         & \multirow{2}{*}{2.5k--32k} & \multirow{2}{*}{40960} \\
                & & Olmo3-Hybrid-7B, Llama3.2-3B & & \\
\bottomrule
\end{tabular}
\caption{\textbf{Per-setting training and evaluation configuration.} Every cell uses the shared
recipe above (full fine-tuning, 1 epoch, lr $5\!\times\!10^{-5}$, global batch 8); only the model
set, the evaluation ladder and the sequence length vary. Ladders are token budgets, evaluated at
2k/4k/8k/16k/32k. Contradiction starts at 2.5k rather than 2k because its realistic ladder begins
at $N{=}56$ --- $N{=}44$ falls below the training minimum of 52. Reordering stops at 16k by ladder
policy. Sequence length is set from each shard's maximum example length, not tuned.}
\label{tab:hyperparams}
\end{claudeblock}
\end{table}

% \begin{tabular}{l l c c c c c}
% \toprule
% \textbf{Task} & \textbf{Model} & $\bm{N}$ & \textbf{seq\_len} & \textbf{mb} & \textbf{ga} & \textbf{epochs} \\
% \midrule
% NQ            & Qwen-3.5-0.8B & $k{=}20$--200 & 4k--32k & 1 & 4 & 4 \\
% HotpotQA      & Qwen-3.5-0.8B & $k{=}20$--200 & 4k--32k & 1 & 4 & 4 \\
% Outlier (wiki, 0.8B) & Qwen-3.5-0.8B & 20--100 & 4k / 16k / 25.6k & 1 & 4 & 1 \\
% Outlier (wiki, 4B)   & Qwen-3.5-4B   & 20--100 & 8k / 16k / 25.6k & 1 & 2 & 1 \\
% Outlier (Amazon)     & Qwen-3.5-0.8B & 30      & 5k       & 1 & 4 & 1 \\
% Grouping             & Qwen-3.5-4B   & 20--100 & 4k / 16k & 1 & 4 & 1 \\
% Contradiction & Qwen-3.5-\{0.8B, 2B, 4B\} & 20--500 & 4k--16k & 1 & 4 & 3 \\
%               & Qwen-2.5-3B               & 20--100 & 4k--16k & 1 & 4 & 3 \\
% Reorder       & Qwen-3.5-9B   & 5--50   & 2k / 4k  & 1 & 1 & 3 \\
% \bottomrule
% \end{tabular}

\paragraph{Mask mixing.}
% ⚠ SUPERSEDED: this described a STATIC mix probability (default alpha=0.10). The project default
% is now a curriculum that anneals p_standard from 0.8 to 0 across training, which is also what
% Section 4.1 of the main text states -- the two contradicted each other.
% Mask-mixing experiments add two YAML keys to the training config:
% \begin{itemize}
% \item \texttt{attention\_pattern: chunked} (or \texttt{random\_token}, \texttt{bigbird}, etc.)
% \item \texttt{standard\_mix\_prob: $\alpha$} -- the probability that a given training step uses the unmasked (full-attention) forward pass instead of the chunked mask. Default $\alpha{=}0.10$ for Qwen-3.5 and $\alpha{=}0.50$ for the older grouping/outlier-scalek runs.
% \item \texttt{mix\_seed: 42} for reproducibility.
% \end{itemize}
\begin{claudeblock}
Mask mixing is a \emph{curriculum} rather than a fixed rate. At each training step a coin with
probability $p_{\text{standard}}$ selects the unmasked (full-attention) forward pass instead of the
document-chunked mask; $p_{\text{standard}}$ anneals linearly from $0.8$ at the first step to $0$
at the last, so training begins close to full attention and ends entirely chunked. The mix seed is
fixed at 42.

One implementation detail is worth stating because it silently voids the curriculum if missed:
the anneal is driven by the number of forward passes, and the data is sharded across data-parallel
ranks, so a per-rank counter leaves $p_{\text{standard}}$ stalled near
$0.8\,(1 - 1/\text{world\_size})$ instead of reaching 0 --- worse the more GPUs are used. Our
launcher predicts the final $p_{\text{standard}}$ before training starts and hard-fails if it misses
0 by more than 0.01.

Evaluation always uses the chunked mask with no mixing, so a mask-mixing result reflects pure
chunked-attention inference and differs from the pure-chunked arm only in how it was trained.
\end{claudeblock}

% ⚠ REMOVED: a "Random-token sparsity" paragraph here referenced Figure~\ref{fig:sparsityscale},
% whose figure and label are no longer in the paper -- one of the two remaining undefined
% references. The random-sparse sweep is not reported in this draft, so the configuration
% paragraph for it has nothing to document. Restore both together if that experiment comes back.
% For the sparsity sweep in Figure~\ref{fig:sparsityscale}, the random-sparse settings use \texttt{attention\_pattern: random\_token} with \texttt{keep\_prob} $\in \{0.10, 0.25, 0.50\}$ and \texttt{random\_seed: 42}. The eval call passes the same \texttt{--keep-prob} and \texttt{--pattern-seed} to materialize the identical random pattern.

\paragraph{Eval defaults.}
% ⚠ SUPERSEDED: eval-set sizes were per-task and several were small (105 for reorder); the suite
% is now a uniform 500 per rung. The backend line is also stale -- both arms now run under vLLM.
% \begin{itemize}
% \item \textbf{Backend:} \texttt{vllm} \cite{kwon2023efficient} for full-attention (faster, FlashAttention \cite{dao2023flashattention2}  kernel) and \texttt{chunked-sdpa} for chunked / random-sparse / anchor variants. NQ, HotpotQA, and outlier full-attention reports use vllm; everything sparse uses chunked-sdpa.
% \item \textbf{Generation:} max\_new\_tokens 200 for contradiction / outlier, 256 for reorder, 1024 for QA tasks. Greedy (temperature 0).
% \item \textbf{Eval set size:} 488 for contradiction, 400 for outlier (wiki), 500 for outlier (Amazon) / NQ / HotpotQA / grouping, 105 for reorder.
% \end{itemize}
\begin{claudeblock}
\begin{itemize}
% ⚠ "500 at every rung" is true of the TOKEN LADDERS only, and must say so: Table~\ref{tab:data-stats}
% reports 488 / 400 / 200 for several of the earlier per-N datasets, and an unqualified claim here
% would contradict it one page earlier. The 488 is the whole contradiction held-out file.
\item \textbf{Eval set size:} 500 examples at every rung of every token ladder, so a difference
below roughly 0.04 is within binomial noise (the standard error is $\pm0.021$ at a score of 0.7 and
$\pm0.010$ at 0.95). The earlier per-$N$ datasets in Table~\ref{tab:data-stats} use the sizes
listed there, of which the smallest is contradiction's 488 --- the entire held-out file.
\item \textbf{Backend:} both arms are served with vLLM \cite{kwon2023efficient}. The chunked arm
applies the document-chunked mask in-process; because that patch cannot reach separate worker
processes, chunked evaluation runs with tensor parallelism 1, and every chunked result is
checked for a non-zero count of masked attention calls before it is recorded. Without that check
a chunked run silently reports the dense number.
\item \textbf{Generation:} greedy (temperature 0). The token budget is set per task from the
rung length plus the maximum answer length; a budget below the prompt length causes the
generation to be skipped and scored zero while the parse rate still reads 1.0, so it is
validated rather than defaulted.
\end{itemize}
\end{claudeblock}

\section{GPU Resources / Compute}
\label{app:gpus}

The experiments for our investigation involved roughly 3k A100 hours and 1k H200 hours (32 CPU cores, 300GB RAM). We expect that with 2xH200 GPUs, the results in Figure~\claude{\ref{fig:multitask_fig4}} would take under 200 hours to reproduce.

% The experiments for this paper were primarily on clusters with SLURM access to 2xA100 and 2xH200 GPU instances. This paper also used VESSL AI cluster resources, with the experiments including a total of roughly 500\$ of credits. For our longer-context experiments (e.g. HPQA 32k context length, contradiction task 4B k=500), in the slower chunked setting, training runs on the H200 cluster could last up to 12-14 hours, with other runs scaling as a function of model scale and context length. The total GPU hour count including experimentation for this project likely totaled around 4k A100 hours. Typically these jobs would be allocated with 32 CPU cores per job, and around 300GB of RAM.

% \amandanote{This is a somewhat unusual amount of detail imo; maybe reframe as (1) total GPU hours (2) we expect that with this type of hardware, it would take X hours to reproduce Y runs?}
% \amandanote{Trying to think about why this feels off to me. I think it's that you can kinda guess that this is from an academic group from these details (SLURM, 2xA100s)?} \prasannnote{Note sure if }